\documentclass[11pt]{article}

\usepackage[utf8]{inputenc}
\usepackage[margin=0.8in]{geometry}
\usepackage{amsmath, amssymb}
\usepackage{booktabs}
\usepackage{xcolor}
\usepackage{enumitem}
\usepackage{hyperref}
\usepackage{tcolorbox} % For highlighted concept boxes
\usepackage{multicol}
\usepackage{array}

% Custom Colors
\definecolor{chemblue}{RGB}{0, 51, 102}
\definecolor{conceptgold}{RGB}{218, 165, 32}

% Formatting
\title{\Huge \textbf{CHM 2045C: Exam 4 Comprehensive Study Guide}}
\author{\Large \textit{Quantum Mechanics, Periodicity, and Chemical Bonding}}
\date{Spring 2026}

\begin{document}

\maketitle

\begin{center}
    \begin{tcolorbox}[colframe=chemblue, colback=white, width=0.8\textwidth, title=Exam Strategy]
        \textbf{Total Questions:} 20 | \textbf{Math:} 3 | \textbf{Conceptual/Numerical:} 17 \\
        This exam focuses heavily on \textit{why} atoms behave the way they do based on electron behavior. Master the trends and Lewis rules to secure the majority of the points.
    \end{tcolorbox}
\end{center}

\section{Chapter 8: The Quantum Mechanical Model}

\subsection{8.1--8.2: Nature of Light}
\textbf{Concept:} Light behaves as both a \textbf{wave} (diffraction/interference) and a \textbf{particle} (photoelectric effect). Particles of light are called \textit{photons}.
\begin{itemize}
    \item \textbf{Wave Formula:} $c = \lambda \nu$ (where $c = 3.00 \times 10^8$ m/s, $\lambda$ is wavelength in meters, $\nu$ is frequency in s$^{-1}$).
    \item \textbf{Particle Formula:} $E = h\nu = \frac{hc}{\lambda}$ (where $h = 6.626 \times 10^{-34}$ J$\cdot$s).
    \item \textbf{Relationship:} Short wavelength ($\lambda$) = High frequency ($\nu$) = High energy ($E$).
\end{itemize}

\subsection{8.4: Wave Properties of Matter (The de Broglie Wavelength)}
\textbf{Concept:} If light can act like a particle, matter (like electrons) can act like a wave.
\begin{tcolorbox}[colframe=chemblue!50, title=Mathematical Method: de Broglie Calculation]
    \textbf{Formula:} $\lambda = \frac{h}{m v}$ \\
    \textbf{Critical Units:} 
    \begin{enumerate}
        \item Mass ($m$) \textbf{MUST} be in \textbf{kilograms} (kg).
        \item Velocity ($v$) \textbf{MUST} be in \textbf{meters per second} (m/s).
        \item $\lambda$ will result in meters (m).
    \end{enumerate}
\end{tcolorbox}

\subsection{8.5: Quantum Numbers}
\textbf{Concept:} Four numbers describe the "address" and properties of an electron in an orbital.
\begin{enumerate}
    \item \textbf{Principal ($n$):} Energy level/Size. $n = 1, 2, 3, \dots$
    \item \textbf{Angular Momentum ($l$):} Shape of the orbital. $l$ ranges from $0$ to $(n-1)$.
    \begin{itemize}
        \item $l=0 \rightarrow s$ (sphere); $l=1 \rightarrow p$ (dumbbell); $l=2 \rightarrow d$ (clover); $l=3 \rightarrow f$.
    \end{itemize}
    \item \textbf{Magnetic ($m_l$):} Orientation in space. $m_l$ ranges from $-l$ to $+l$ (including 0).
    \item \textbf{Spin ($m_s$):} Direction of electron spin. Either $+1/2$ or $-1/2$.
\end{enumerate}

\subsection{8.5: Hydrogen Electron Transitions (Bohr Model)}
\textbf{Mathematical Method:} Calculating energy change when an electron jumps levels.
\[ \Delta E = -2.18 \times 10^{-18} \text{ J} \left( \frac{1}{n_f^2} - \frac{1}{n_i^2} \right) \]
\begin{itemize}
    \item \textbf{Absorption:} $n_f > n_i$ ($\Delta E$ is positive).
    \item \textbf{Emission:} $n_f < n_i$ ($\Delta E$ is negative; light is released).
\end{itemize}

\newpage

\section{Chapter 9: Periodic Properties}

\subsection{9.3: Rules for Electron Configuration}
\begin{itemize}
    \item \textbf{Pauli Exclusion Principle:} No two electrons in an atom can have the same four quantum numbers (max 2 electrons per orbital, must have opposite spins).
    \item \textbf{Hund's Rule:} When filling degenerate orbitals (same energy), electrons fill them singly first with parallel spins to minimize repulsion.
    \item \textbf{Aufbau Principle:} Electrons fill the lowest energy orbitals first ($1s \rightarrow 2s \rightarrow 2p \dots$).
\end{itemize}

\subsection{9.6--9.8: Periodic Trends}
\textbf{Underlying Concept:} Trends are driven by \textbf{Effective Nuclear Charge ($Z_{\text{eff}}$)}, which is the net positive charge "felt" by valence electrons after accounting for shielding by core electrons.

\begin{tcolorbox}[colframe=conceptgold, title=Visualizing Trends (Arrows point toward INCREASE)]
    \begin{center}
        \begin{tabular}{l l l}
            \textbf{Trend} & \textbf{Horizontal ($\leftrightarrow$)} & \textbf{Vertical ($\updownarrow$)} \\ \midrule
            Atomic Radius & $\longleftarrow$ (Increase Left) & $\downarrow$ (Increase Down) \\
            Ionization Energy (IE) & $\longrightarrow$ (Increase Right) & $\uparrow$ (Increase Up) \\
            Metallic Character & $\longleftarrow$ (Increase Left) & $\downarrow$ (Increase Down) \\
            Electron Affinity & $\longrightarrow$ (More Negative Right) & $\uparrow$ (More Negative Up) \\
            Effective Nuclear Charge & $\longrightarrow$ (Increase Right) & (Slight Increase Down) \\ \bottomrule
        \end{tabular}
    \end{center}
\end{tcolorbox}

\subsection{9.7: Magnetic Properties}
\begin{itemize}
    \item \textbf{Paramagnetic:} Atom or ion has \textbf{unpaired} electrons. Attracted to magnetic fields.
    \item \textbf{Diamagnetic:} All electrons are \textbf{paired}. Slightly repelled by magnetic fields.
\end{itemize}

\section{Chapter 10: Chemical Bonding I}

\subsection{10.5--10.6: Bond Classification}
Determined by the difference in electronegativity ($\Delta\text{EN}$):
\begin{itemize}
    \item \textbf{0.0 -- 0.4:} Pure (Nonpolar) Covalent (Equal sharing).
    \item \textbf{0.4 -- 2.0:} Polar Covalent (Unequal sharing).
    \item \textbf{> 2.0:} Ionic (Electron transfer).
\end{itemize}

\subsection{10.8: Formal Charge}
Used to identify the \textbf{best} Lewis structure among resonance forms.
\textbf{Formula:} $\text{FC} = (\text{Valence } e^-) - [\text{Unpaired } e^- + \frac{1}{2}(\text{Bonding } e^-)]$
\textit{Selection Rules:}
1. Smaller FCs are better (closer to zero).
2. Negative FCs should be on the more electronegative atom.

\subsection{10.10: Calculating Bond Energy ($\Delta H_{\text{rxn}}$)}
\begin{tcolorbox}[colframe=chemblue, title=Meticulous Method: Step-by-Step $\Delta H_{\text{rxn}}$]
    \textbf{Formula:} $\Delta H_{\text{rxn}} = \sum (\text{Bonds Broken}) - \sum (\text{Bonds Formed})$
    \begin{enumerate}
        \item \textbf{Draw:} Sketch the Lewis structures for \textbf{every} reactant and product.
        \item \textbf{Inventory Reactants (Breaking):} Count every bond in the reactants. Multiply by the coefficients from the balanced equation. These values are \textbf{Positive}.
        \item \textbf{Inventory Products (Forming):} Count every bond in the products. Multiply by the coefficients. These values are \textbf{Negative} in the formula.
        \item \textbf{Calculate:} [Sum of Reactant Bond Energies] $-$ [Sum of Product Bond Energies].
    \end{enumerate}
    \textit{Note: Breaking bonds costs energy (+), forming bonds releases energy (-).}
\end{tcolorbox}

\newpage

\section{Chapter 11: Molecular Shapes (VSEPR)}

\subsection{11.2--11.4: VSEPR Theory}
\textbf{Concept:} Electron groups (lone pairs and bonds) repel each other and want to be as far apart as possible.
\begin{itemize}
    \item \textbf{Electron Geometry:} Based on total electron groups.
    \item \textbf{Molecular Geometry:} Based only on where atoms are (lone pairs change the name).
\end{itemize}

\begin{center}
\renewcommand{\arraystretch}{1.5}
\begin{tabular}{|c|c|c|l|c|}
\hline
\textbf{Total Groups} & \textbf{Lone Pairs} & \textbf{Electron Geo} & \textbf{Molecular Geo} & \textbf{Ideal Angle} \\ \hline
2 & 0 & Linear & Linear & $180^\circ$ \\ \hline
3 & 0 & Trigonal Planar & Trigonal Planar & $120^\circ$ \\ 
3 & 1 & Trigonal Planar & \textbf{Bent} & $<120^\circ$ \\ \hline
4 & 0 & Tetrahedral & Tetrahedral & $109.5^\circ$ \\ 
4 & 1 & Tetrahedral & \textbf{Trigonal Pyramidal} & $<109.5^\circ$ \\ 
4 & 2 & Tetrahedral & \textbf{Bent} & $<<109.5^\circ$ \\ \hline
5 & 0 & Trig. Bipyramidal & Trig. Bipyramidal & $90^\circ, 120^\circ$ \\ \hline
6 & 0 & Octahedral & Octahedral & $90^\circ$ \\ \hline
\end{tabular}
\end{center}

\subsection{11.3: Lone Pair Effects}
\textbf{Concept:} Lone pairs are "physically larger" than bonding pairs because they are held by only one nucleus. They exert more repulsion, which \textbf{squeezes} the bonding angles to be smaller than the ideal values.

\section{Summary Checklist}
\begin{itemize}
    \item [ ] Can I convert units for de Broglie (g $\rightarrow$ kg)?
    \item [ ] Do I remember that $l$ values $0, 1, 2, 3$ correspond to $s, p, d, f$?
    \item [ ] Can I rank atoms by size: Fr (largest) vs F (smallest)?
    \item [ ] Do I know that $Z_{\text{eff}}$ increases to the \textbf{right}?
    \item [ ] Can I calculate Formal Charge for a molecule like $N_2O$?
    \item [ ] Am I remembering to subtract "Formed" from "Broken" for $\Delta H$?
\end{itemize}

\end{document}
